\documentclass[11pt]{article}

\usepackage[utf8]{inputenc}
\usepackage{graphicx}
\usepackage[dvipsnames]{xcolor}
\usepackage{color}
\usepackage[margin=1in]{geometry}
\usepackage{hyperref}
\usepackage{tikz}
\usepackage{amsmath}
\usepackage{commath}
\usepackage{relsize}
\usepackage{centernot}
\usepackage{amssymb}
\usepackage{amsfonts}
\usepackage{graphicx}
\usepackage{textcomp}
\usepackage{gensymb}
\usepackage{enumitem}
\usepackage{siunitx}
\usepackage{fancyhdr}
\usepackage{floatrow}
\usepackage{wrapfig}
\usepackage[labelfont=bf]{caption}
\usepackage{mathtools}
\usepackage{pgfplots}
\usepackage{multicol}
\usepackage[misc]{ifsym}
\usepackage{tabularx}
\usepackage{empheq}
\usepackage{tkz-tab}
\usepackage{xpatch}
\usepackage{caption}
\usepackage{subcaption}
\usepackage{tabularray}
\usepackage{esint}
\usepackage{amsthm}
\usepackage{thmtools}
\RequirePackage[most]{tcolorbox}  

\swapnumbers

\makeatletter
\declaretheoremstyle[
  headfont= \sffamily\bfseries\color{cyan!50!black},
  headpunct={\sffamily\bfseries.},
  postheadspace=0.5em,
  notefont=\sffamily\bfseries,
  headformat=\NAME\NUMBER\let\thmt@space\@empty\NOTE,
  bodyfont=\mdseries\slshape,
  spaceabove=12pt,spacebelow=12pt,
]{thmstyle}
\makeatother

\theoremstyle{thmstyle}
\declaretheorem[name=Definición]{theorem}
\newtheorem{definition}[theorem]{Definición}

\definecolor{lightcyan}{rgb}{0.88, 1.0, 1.0}
\colorlet{mythmback}{lightcyan!40!white}

\tcolorboxenvironment{theorem}{
  colback=mythmback,coltitle=blue,colframe=mythmback,
  left=0pt,right=0pt,
  top=0pt,bottom=0pt,
  before skip=10pt, after skip=10pt,
}

\xpatchcmd{\tkzTabLine}{$0$}{$\bullet$}{}{}
\tikzset{t style/.style={style=solid}}
\pgfplotsset{compat=newest}

\newtheorem{manualtheoreminner}{}
\newenvironment{manualtheorem}[1]{%
  \renewcommand\themanualtheoreminner{#1}%
  \manualtheoreminner
}{\endmanualtheoreminner}

\newcommand{\pd}[2]{\frac{\partial #1}{ \partial #2}}
\newcommand{\td}[2]{\frac{\mathrm{d} #1}{ \mathrm{d} #2}}
\newcommand{\pdd}[2]{\frac{\partial^2 #1}{ \partial #2 ^2}}
\newcommand{\pddc}[3]{\frac{\partial^2 #1}{ \partial #2 \partial #3}}
\newcommand{\solution}[0]{\large{\textbf{Solución}}\normalsize}

\def\sca   #1{\mbox{\rm{#1}}{}}
\def\mat   #1{\mbox{\boldmath $\mathsf #1$}}
\def\vec   #1{\mbox{\boldmath $#1$}{}}
\def\ten   #1{\mbox{\boldmath $#1$}{}}
\def\ltr   #1{\mbox{\sf{#1}}}
\def\bltr  #1{\mbox{\sffamily{\bfseries{{#1}}}}}

\DeclareMathOperator{\dive}{div}
\DeclareMathOperator{\trace}{trace}
\DeclareMathOperator{\tr}{tr}
\DeclareMathOperator{\symm}{symm}
\DeclareMathOperator{\sk}{skew}
\DeclareMathOperator{\grad}{grad}
\DeclareMathOperator{\Grad}{Grad}
\DeclareMathOperator{\curl}{curl}
\DeclareMathOperator{\Curl}{Curl}
\DeclareMathOperator*{\argmin}{\arg\!\min}

\def\dx{\mbox{\(\,\mathrm{d}x\)}}
\def\ds{\mbox{\(\,\mathrm{d}s\)}}
\def\dS{\mbox{\(\,\mathrm{d}S\)}}
\def\dV{\mbox{\(\,\mathrm{d}V\)}}
\def\dv{\mbox{\(\,\mathrm{d}v\)}}
\def\R{\mathbb R}
\def\C{\mathbb C}
\def\N{\mathbb N}
\def\Q{\mathbb Q}
\def\Z{\mathbb Z}

\usepackage{mdframed}
\mdfdefinestyle{MyFrame}{%
    linecolor=black,
    linewidth=1pt,
    outerlinewidth=1pt,
    roundcorner=20pt,
    innertopmargin=\baselineskip,
    innerbottommargin=\baselineskip,
    innerrightmargin=20pt,
    innerleftmargin=20pt,
    splittopskip=2\baselineskip,
    backgroundcolor=gray!50!white}
    
\setlength\textwidth{7in}
\setlength\parindent{0pt}
\oddsidemargin=-0.5cm
\pagestyle{fancy}
\newcommand*{\vertbar}{\rule[-1ex]{0.5pt}{2.5ex}}
\setlength{\headheight}{13.59999pt}

\renewcommand{\headrulewidth}{0.1pt}
\renewcommand{\footrulewidth}{0.1pt}
\renewcommand{\figurename}{Figura}
\renewcommand\refname{Referencias}
\renewcommand{\thesection}{\arabic{section}.}
\renewcommand{\thesubsection}{\thesection\arabic{subsection}.}
\renewcommand{\labelitemi}{{\raisebox{.3\height}{\scalebox{0.6}{$\blacklozenge$}}}}
\renewcommand*\contentsname{Índice}

\newcommand{\p}{\vspace{10pt}}

\AtBeginDocument{\vspace*{-3.2\baselineskip}}
\textheight=665pt

\lhead{\scshape Pontificia Universidad Católica de Chile}
\rhead{\scshape IMC UC}
\cfoot{\thepage}

\usepackage{footnote}
\usepackage{etoolbox}
\newtoggle{indefintion}
\togglefalse{indefintion}
\pretocmd{\footnote}{\iftoggle{indefintion}{\stepcounter{footnote}}{\relax}}{}{}

\begin{document}
    \begin{flushleft}
        \small
        \vspace{1pt}
        \textbf{IMT3130} \hfill
        \textbf{20/03/2026}
    \end{flushleft}
    \begin{center}
        \LARGE\textbf{Ayudantía 2}\\
        \vspace{3pt}
        \normalsize\textbf{Diferencias finitas}\\
        \normalsize Ayudante: Bastián Herrera \Letter{\hspace{1pt}: \texttt{biherrera@uc.cl}}
        \rule{\textwidth}{0.05mm}
    \end{center}

\section{Problemas}

\begin{enumerate}[label=\textbf{\arabic*.}, leftmargin=*]

    \item Sea $f$ suficientemente suave y considere la familia de aproximaciones
    \[
    \frac{a f(x_i-2h)+b f(x_i-h)+c f(x_i+h)+d f(x_i+2h)}{h},
    \]
    donde $a,b,c,d\in\R$ son constantes independientes de $h$.
    Determine $a,b,c,d$ de modo que la expresión aproxime $f'(x_i)$ con el mayor orden posible. Escriba el esquema resultante y determine su orden de consistencia.

    \item Considere el problema de borde
    \[
    -u''(x)=f(x), \qquad x\in(0,1), \qquad u(0)=u(1)=0,
    \]
    y su discretización estándar
    \[
    \frac{-u_{i+1}+2u_i-u_{i-1}}{h^2}=f(x_i), \qquad i=1,\dots,N.
    \]
    Suponga además que el error de truncamiento local es de orden $O(h^2)$. ¿Esto basta para concluir que el error global también es de orden $O(h^2)$? Deduzca la ecuación de error, identifique la hipótesis adicional necesaria para obtener convergencia, y discuta qué ocurre si solo se dispone de una cota de estabilidad de la forma
    \[
    \|v_h\|\le C_h\|A_hv_h\|,
    \qquad C_h\to\infty \text{ cuando } h\to 0.
    \]

    \item Sea $\Omega=(0,1)^2$ y considere una malla cartesiana uniforme
    \[
    x_i=ih,\qquad y_j=jh.
    \]
    Para nodos interiores, sea $u_{ij}\approx u(x_i,y_j)$. Construya un esquema de diferencias finitas para
    \[
    -\Delta u=f
    \qquad \text{en } \Omega,
    \]
    usando diferencias centradas en cada dirección. Escriba el \textit{stencil} resultante y pruebe, mediante expansiones de Taylor en dos variables, que el método es consistente de orden $2$. Identifique además el término dominante del error local.

\end{enumerate}

\end{document}